%% LyX 2.4.2.1 created this file.  For more info, see https://www.lyx.org/.
%% Do not edit unless you really know what you are doing.
\documentclass[acmlarge]{acmart}
\usepackage[T1]{fontenc}
\usepackage[utf8]{inputenc}
\usepackage{graphicx}
\usepackage{babel}
\usepackage{subcaption} % Added for better side-by-side figures
\usepackage{placeins}



\title{Machine Learning Project: DDoS Attacks Detection and Characterization}


\author{Angélica Castillo}
\affiliation{
  \institution{Politecnico di Torino}
  \city{Turin}
  \country{Italy}
}
\email{s354702@studenti.polito.it}

\author{Roey Bretschneider}
\affiliation{
  \institution{Politecnico di Torino}
  \city{Turin}
  \country{Italy}
}
\email{s350457@studenti.polito.it}

\author{Antonio Manconi}
\affiliation{
  \institution{Politecnico di Torino}
  \city{Turin}
  \country{Italy}
}
\email{s355115@studenti.polito.it}


\begin{abstract}
This project investigates the use of machine learning techniques for the analysis and detection of Distributed Denial of Service (DDoS) attacks in network traffic. Starting from flow-level data extracted using CICFlowMeter, an extensive exploratory analysis and feature engineering process is conducted to characterize benign and malicious traffic patterns. The resulting feature matrix is then leveraged for supervised and unsupervised learning tasks aimed at identifying and classifying DDoS attacks. The results highlight the importance of combining volumetric, temporal, and protocol-level features to effectively model complex network traffic behaviors.
\end{abstract}

\begin{CCSXML}
<ccs2012>
   <concept>
       <concept_id>10010147.10010257</concept_id>
       <concept_desc>Computing methodologies~Machine learning</concept_desc>
       <concept_significance>500</concept_significance>
   </concept>
   <concept>
       <concept_id>10003033.10003079</concept_id>
       <concept_desc>Networks~Network security</concept_desc>
       <concept_significance>500</concept_significance>
   </concept>
</ccs2012>
\end{CCSXML}

\ccsdesc[500]{Computing methodologies~Machine learning}
\ccsdesc[500]{Networks~Network security}

\keywords{DDoS detection, network traffic analysis, machine learning, network security, flow-based features, traffic classification.}

\settopmatter{printacmref=false}
\setcopyright{none}
\renewcommand\footnotetextcopyrightpermission[1]{}

\setlength{\textfloatsep}{10pt}    % espacio entre figura y texto
\setlength{\floatsep}{8pt}         % espacio entre dos figuras
\setlength{\intextsep}{10pt}       % figuras dentro del texto

\begin{document}
\maketitle
\clearpage

\section{Data Exploration and Pre-processing}

\subsection{Dataset Overview}

The dataset used in this project contains network traffic flows extracted using CICFlowMeter-V3, where each flow represents a bidirectional communication between a source and a destination IP within a specific time window. Each flow is characterized by a set of statistical features computed from packet-level information, including temporal, volumetric, and protocol-related attributes.

The dataset includes both benign traffic and multiple types of modern Distributed Denial of Service (DDoS) attacks. Each flow is associated with a ground truth (GT) label that identifies the corresponding traffic type, enabling supervised analysis and performance evaluation in subsequent tasks.

Overall, the dataset aims to realistically reproduce real-world network conditions by combining benign user-generated traffic with multiple attack scenarios, making it suitable for traffic analysis, anomaly detection, and attack characterization.


\subsection{Generic Traffic-Level Analysis}

As a first step, an exploratory analysis was conducted at a generic traffic level, independently of the ground truth labels, in order to understand the overall behavior and statistical properties of the network flows.

\begin{figure}[H]
    \centering
    \begin{subfigure}{0.48\linewidth}
        \centering
        \includegraphics[width=\linewidth]{Flow_packets.png}
        \caption{Distribution of Flow Packets per second at the generic traffic level.}
    \end{subfigure}
    \hfill
    \begin{subfigure}{0.48\linewidth}
        \centering
        \includegraphics[width=\linewidth]{Flow_byte.png}
        \caption{Distribution of Flow Bytes per second highlighting traffic load variability.}
    \end{subfigure}
\end{figure}

The previous figure shows the distribution of the number of packets per second generated by each flow. The distribution exhibits a heavy-tailed behavior, typical of large-scale network traffic and potentially indicative of anomalous activity.

Similarly, the distribution of flow bytes per second highlights a wide variability in traffic load across flows. A limited number of flows contribute disproportionately to the total traffic volume, suggesting the coexistence of low-rate benign traffic and high-rate traffic bursts potentially associated with DDoS attacks.


\subsection{Ground Truth (GT) Level Analysis}

After the initial exploration at a generic traffic level, the analysis was extended by incorporating the ground truth (GT) labels in order to investigate differences between benign traffic and various DDoS attack types.

Figure 3 illustrates the relationship between Flow Packets per second and Flow Bytes per second using a logarithmic scale, with flows colored according to their GT labels. The visualization reveals a clear separation trend between benign traffic and attack traffic. Benign flows are mainly concentrated in regions characterized by lower packet and byte rates, while attack flows tend to occupy areas with significantly higher traffic intensity.

In addition, boxplot analyses were performed to compare the distributions of selected features across different traffic types. Figure 4 shows the distribution of Flow Duration for benign and attack traffic.

\begin{figure}[H]
    \centering
    \begin{subfigure}{0.48\linewidth}
        \centering
        \includegraphics[width=\linewidth]{GT_1.png}
        \caption{Relationship between Flow Packets per second and Flow Bytes per second in logarithmic scale, with flows colored according to ground truth labels}
    \end{subfigure}
    \hfill
    \begin{subfigure}{0.48\linewidth}
        \centering
        \includegraphics[width=\linewidth]{GT_2.png}
        \caption{Distribution of Flow Duration for benign traffic and different DDoS attack types}
    \end{subfigure}
\end{figure}

Overall, the GT-level analysis demonstrates that significant differences exist between benign traffic and various types of DDoS attacks. These differences are clearly observable through data visualization techniques and motivate the use of machine learning methods to automatically detect and classify malicious traffic.


\subsection{Feature Engineering}

Based on the insights obtained from the exploratory analysis, additional features were generated to better capture the underlying traffic behavior and enhance the discriminative power of the dataset for subsequent machine learning tasks.

First, volumetric features were derived to quantify the overall traffic associated with each flow. In particular, the total number of packets per flow was computed as the sum of forward and backward packets, while the total number of bytes was obtained by aggregating the total length of forward and backward packets. These features provide a compact representation of the traffic volume exchanged during each communication flow.

In addition to the original flow-level features provided by CICFlowMeter, a set of derived features was introduced to better capture structural, directional, and temporal characteristics of network traffic flows. The motivation behind this feature engineering process was to enrich the representation of each flow with indicators that are less dependent on absolute traffic volume and more focused on behavioral patterns, which are particularly relevant in the context of DDoS attack detection and characterization.

Specifically, the engineered features were designed to capture complementary aspects of traffic behavior, with particular emphasis on directionality and flow asymmetry. Directionality-related features quantify the imbalance between forward and backward traffic, highlighting flows that exhibit strong asymmetric or unidirectional behavior, which is a common trait of many flooding and reflection-based attacks. Ratio-based features, expressed as normalized fractions rather than absolute counts, were employed to describe how packets, bytes, and protocol flags are distributed across directions within a flow, making these features more robust to scale variations.

Moreover, variability-oriented features were introduced to characterize the temporal and structural regularity of traffic. Coefficients of variation derived from inter-arrival times and packet lengths provide insight into whether a flow exhibits highly regular, machine-generated behavior or more heterogeneous dynamics typically associated with benign application traffic. Finally, aggregated volumetric features summarizing the total number of packets and bytes exchanged within a flow were included to retain information about overall traffic intensity, while still allowing this information to be interpreted jointly with directional and temporal indicators.

Overall, these engineered features aim to complement the original dataset by providing a richer and more expressive description of network traffic behavior. By jointly modeling volume, directionality, temporal regularity, and protocol-level characteristics, the resulting feature set enhances the discriminative power of the dataset and supports both supervised and unsupervised learning tasks in subsequent sections.

\subsection{Data Pre-processing}

Prior to applying machine learning algorithms, a data pre-processing phase was performed to improve the quality and suitability of the feature matrix. This phase focused on handling feature scaling, analyzing feature correlations, and investigating the intrinsic structure of the data through dimensionality reduction techniques.

\subsubsection{Correlation Analysis}

A correlation analysis was conducted on the numerical features to identify redundant information and strong linear dependencies. The correlation matrix revealed the presence of multiple feature pairs exhibiting very high absolute correlation values, indicating potential multicollinearity issues.

To mitigate this effect, features with an absolute correlation greater than a predefined threshold were removed from the dataset. This step reduces redundancy, simplifies the feature space, and helps improve the stability and interpretability of subsequent machine learning models, without significantly affecting the information content of the dataset.

\begin{figure}[H]
    \centering
    \includegraphics[width=0.5\linewidth]{Heatmap.png}
    \caption {Correlation heatmap of numerical features highlighting strong linear dependencies among several feature pairs}
\end{figure}

\subsubsection{Feature Scaling}

Given the heterogeneous nature of the extracted features, which span different ranges and physical units, feature scaling was required. Standardization was applied to transform each feature to have zero mean and unit variance. This step is particularly important for distance-based methods and dimensionality reduction techniques, such as Principal Component Analysis (PCA), as well as for many machine learning algorithms used in the subsequent sections.

\subsubsection{Principal Component Analysis}

Principal Component Analysis was applied to the standardized feature matrix as an exploratory tool to visualize the data in a reduced-dimensional space. By projecting the high-dimensional data onto the first two principal components, it was possible to analyze the global structure of the dataset and observe how different traffic types are distributed in the feature space.

The PCA visualization shows that some degree of clustering exists among flows belonging to the same attack type, while partial overlaps between different classes are also observed. This behavior highlights the intrinsic complexity of DDoS traffic patterns and suggests that linear separability is limited, thus motivating the adoption of more expressive machine learning models in later stages of the project.

Overall, the pre-processing steps performed in this section aim to enhance data quality, reduce redundancy, and provide a clearer understanding of the dataset structure, thereby preparing a reliable and well-conditioned feature matrix for the supervised and unsupervised learning tasks presented in the following sections.

\begin{figure}[H]
    \centering
    \includegraphics[width=0.4\linewidth]{PCA.png}
    \caption {PCA projection of standardized features onto the first two principal components. The percentage of explained variance is reported for each component, and flows are colored according to ground truth labels}
\end{figure}

\subsection{Dataset Construction Strategy}

To properly address both supervised and unsupervised learning tasks while avoiding methodological pitfalls, a specific dataset construction strategy was adopted, with particular attention paid to preventing information leakage during feature selection and dimensionality reduction.

Correlation analysis and dimensionality reduction techniques, such as Principal Component Analysis (PCA), are commonly used to mitigate feature redundancy and multicollinearity. However, when applied before splitting the data into training and test sets, these techniques may introduce information leakage, as statistical relationships are computed using samples that should remain unseen during model training. This issue is especially critical in supervised learning scenarios, where reliable model generalization must be ensured.

For this reason, two different feature matrices were generated, each tailored to the requirements of the subsequent learning tasks. For the supervised classification task, the dataset was first split into training and test sets, and correlation analysis, feature selection, and dimensionality reduction were performed exclusively on the training data. The resulting transformations were subsequently applied to the test set, preserving a strict separation between training and evaluation data.

In contrast, for the unsupervised clustering task, no ground truth labels are used during model training or evaluation. Consequently, preprocessing was performed on the entire dataset after removing label-related information, allowing the algorithms to exploit the global structure of the data and capture intrinsic similarities among traffic flows.

Overall, this dual-dataset strategy ensures methodological correctness while enabling each learning task to benefit from an appropriate preprocessing pipeline, achieving reliable performance evaluation and meaningful traffic characterization.

\subsection{Final Feature Matrix Characterization}

Overall, the final feature matrix provides a compact yet expressive representation of network traffic behavior. The observed distributional differences between benign and malicious traffic, together with the reduced correlation among features, suggest that the dataset is well-prepared for both supervised classification and unsupervised clustering tasks. This concludes the data exploration and pre-processing phase and establishes a solid foundation for the machine learning models presented in the following sections.


\section{Supervised learning -- classification}


\subsection{Description of Splits}
To rigorously evaluate generalization, we employed two distinct stratified splitting strategies.
\\
\textbf{Per-Flow Split (Stratified Random):} Flows were shuffled randomly and assigned to training or testing sets independent of time. This serves as a theoretical baseline assuming no temporal dependency between flows.
\\
\textbf{Per-Time Split (Stratified Temporal):} Samples were sorted chronologically. The model is trained on past data and tested on future data to evaluate robustness against concept drift and changing attack signatures.

\subsection{``On the fly'' processing}
To prevent data leakage, preprocessing steps were integrated directly into the training pipeline. Specific identifiers (IPs, Timestamp) were removed to force learning of generalizable behaviors. Feature scaling (Standardization) and dimensionality reduction were fit dynamically on the training fold and applied to the test fold to ensure the model does not ``peek'' at the test distribution.

\subsection{Initial Model Screening}
We established a baseline using Logistic Regression (Linear), Gaussian Naive Bayes (Probabilistic), Decision Tree, and Random Forest (Ensemble) with default hyperparameters.
\subsubsection{Baseline Performance}
Results indicate that Logistic Regression and GNB underfit the data, failing to capture complex non-linear traffic patterns. Conversely, tree-based models achieved higher accuracy with only mild overfitting. Decision Tree and Random Forest yielded the best initial performance across both splits.

\begin{figure}[H]
    \centering
    \begin{subfigure}[b]{0.48\linewidth}
        \includegraphics[width=\linewidth]{base_models_stratified_bar}
        \caption{Stratified Split}
    \end{subfigure}
    \hfill
    \begin{subfigure}[b]{0.48\linewidth}
        \includegraphics[width=\linewidth]{base_models_stratified_time_bar}
        \caption{Temporal Split}
    \end{subfigure}
    \caption{Baseline Model Performance}
\end{figure}

\subsubsection{results - confusion matrix}

\leavevmode

% --- Plot 1: Stratified Split ---
\begin{figure}[H]
    \centering
    % Subfigure wrapper removed as it is now a standalone figure
    \includegraphics[width=0.9\linewidth]{base_models_stratified_cm}
    % Combined the previous main caption and sub-caption
    \caption{Baseline Models: Stratified Split Confusion Matrix}
    % Don't forget the description for acmart!
    \Description{Confusion matrix for the stratified split showing...} 
    \label{fig:stratified_cm}
\end{figure}

% --- Plot 2: Temporal Split ---
\begin{figure}[H]
    \centering
    \includegraphics[width=0.9\linewidth]{base_models_stratified_time_cm}
    \caption{Baseline Models: Temporal Split Confusion Matrix}
    \Description{Confusion matrix for the temporal split showing...}
    \label{fig:temporal_cm}
\end{figure}


Two consistent patterns emerged: (1) \textbf{Class Overlap:} Class 3 is systematically misclassified as Class 2, and Classes 10/11 exhibit bidirectional confusion, suggesting these pairs share indistinguishable feature signatures. (2) \textbf{Concept Drift in GNB:} While tree models and LR remained robust, GNB failed completely on Class 7 in the Temporal Split (misclassifying it as Class 2). This confirms that Class 7's statistical distribution shifted over time, breaking GNB's rigid assumptions.

\subsubsection{classification report}

we can see that our Decision Tree baseline model perform very good
on benign label and has bad performance on label 3,10,11 (we do not
elaborate more here due to limitation of pages on our work)


\begin{figure}[H]
    \centering
    \includegraphics[width=0.4\linewidth]{classification_report_DT_baseline}
    \caption{Classification Report for Decision Tree Baseline}
\end{figure}


\subsection{Hyperparameter Tuning}
We implemented a hierarchical tuning strategy using Cross-Validation, conducted independently for each split to prevent leakage. We first performed a coarse sensitivity analysis (Impact Tuning) to identify critical hyperparameters, followed by a focused 2D Grid Search for optimization.

\subsubsection{impact tuning}
we will present here the results only for the first split due to limitation
of pages for our work.\\
\leavevmode


\begin{figure}[H]
    \centering
    \includegraphics[width=0.8\linewidth]{decisiontreeclassifier_params_impact_stratified_split_impact}
    \caption{Decision Tree Parameter Impact}
\end{figure}

\begin{figure}[H]
    \centering
    \begin{subfigure}[b]{0.48\linewidth}
        \includegraphics[width=\linewidth]{randomforestclassifier_params_impact_stratified_split_impact}
        \caption{Random Forest}
    \end{subfigure}
    \hfill
    \begin{subfigure}[b]{0.48\linewidth}
        \includegraphics[width=\linewidth]{logisticregression_params_impact_stratified_split_impact}
        \caption{Logistic Regression}
    \end{subfigure}
    \caption{Hyperparameter Impact Analysis}
\end{figure}

according to those results we decided that the best hyperparameters
for logistic regression or C= None, panelty = None, solver=lbfgs.

for Random Forest Classifier we chose to do 2d scan on max depth and
min\_samples\_split while taking min\_samples\_leaf to be 1 due to
his better results and n\_estimators to 100 because we saw it didn't
change much and we wanted to save runtime. for Decision Tree we chose
to do 2d scan on max depth and criterion while taking min\_samples\_leaf
to be 1 due to his better results, min\_samples\_split=5, max\_features=None.

\subsubsection{2d grid search}

\leavevmode

\begin{figure}[H]
    \centering
    \begin{subfigure}[b]{0.48\linewidth}
        \includegraphics[width=\linewidth]{decisiontreeclassifier_2d_grid_search_stratified_split_min_samples_split5_min_samples_leaf1_max_featuresnone_grid2d}
        \caption{Decision Tree}
    \end{subfigure}
    \hfill
    \begin{subfigure}[b]{0.48\linewidth}
        \includegraphics[width=\linewidth]{randomforestclassifier_2d_grid_search_stratified_split_n_estimators100_min_samples_leaf1_grid2d}
        \caption{Random Forest}
    \end{subfigure}
    \caption{2D Grid Search Results}
\end{figure}

For the Random Forest we took \texttt{max\_depth} of \texttt{None} with \texttt{min\_samples\_split}=14 which yielded the best results of 0.78. \\
For the Decision Tree we took \texttt{max\_depth}=16 and \texttt{criterion} of \texttt{gini} which yielded the best results of 0.7765.

\subsubsection{validation curve}

for the GNB model there was only one hyperparameter to tune so we
used the validation curve in order to take the best configuration
for it. 

\begin{figure}[H]
    \centering
    \includegraphics[width=0.3\linewidth]{stratified_split_var_smoothing_val_curve}
    \caption{Validation Curve for GNB}
\end{figure}


\subsection{Efficiency \& Timing Analysis}
We leveraged multi-core processing to accelerate training (which had significant impact on RF results due to training of multipe trees at parallel ). As expected, \textbf{Inference Time} for the Random Forest was significantly higher than the Decision Tree due to the ensemble overhead (100 estimators vs. 1 tree). Regarding \textbf{Training Time}, Logistic Regression proved computationally expensive as we increased the maximum iterations ($max\_iter=3000$) to ensure convergence.

\subsubsection{results}

\leavevmode

\begin{figure}[H]
    \centering
    \begin{subfigure}[b]{0.48\linewidth}
        \includegraphics[width=\linewidth]{inference_time_comparison_bar}
        \caption{Inference Time}
    \end{subfigure}
    \hfill
    \begin{subfigure}[b]{0.48\linewidth}
        \includegraphics[width=\linewidth]{train_time_comparison_bar}
        \caption{Training Time}
    \end{subfigure}
    \caption{Performance Timing Comparison}
\end{figure}

\subsection{final results}

\subsubsection{results after tuning}

\leavevmode

\begin{figure}[H]
    \centering
    \begin{subfigure}[b]{0.48\linewidth}
        \includegraphics[width=\linewidth]{tuned_model_stratified_split_bar}
        \caption{Stratified Split}
    \end{subfigure}
    \hfill
    \begin{subfigure}[b]{0.48\linewidth}
        \includegraphics[width=\linewidth]{tuned_model_temporal_stratified_split_bar}
        \caption{Temporal Split}
    \end{subfigure}
    \caption{Tuned Model Performance}
\end{figure}

\begin{figure}[H]
    \centering
    \includegraphics[width=0.9\linewidth]{tuned_model_stratified_split_cm}
    \caption{Tuned Model Confusion Matrix (Stratified)}
\end{figure}

\begin{figure}[htbp]
    \centering
    \includegraphics[width=0.9\linewidth]{tuned_model_temporal_stratified_split_cm}
    \caption{Tuned Model Confusion Matrix (Temporal)}
\end{figure}

\subsubsection{Analysis of Tuning Results}

It is important to note that while our hyperparameter tuning process
maximized the Cross-Validation accuracy on the training set, it did
not strictly guarantee improved performance on the unseen Test Set.
In some instances, the tuned models performed marginally lower than
the baseline. this happened due to the fact that our model chose the
best results for different validation folds, In contrast, the library
default parameters are often designed to be \textquotedbl conservative\textquotedbl{}
and robust across varied distributions. Consequently, when the \textquotedbl optimized\textquotedbl{}
model was applied to the Test Set---which contains unseen attack
patterns --- its specialized configuration generalized slightly less
effectively than the more generic baseline model.\\
it seems most of the results didn't change significantly but the GNB
model on the temporal stratified split, if we take a closer look at
the relevant confusion matrix we can notice that our false prediction
of class 2 instead of class 7 reduce significantly

\subsection{results investigation}

\subsubsection{model selection}

based on the training results our best model is the Decision Tree.
it supplied the best results with the RF but beats him significantly
in aspect of inference and training time. it's also worth mentioning
that all models were pretty good distinguishing class 0 which is the
benign in comparison with other attacks.

\subsubsection{FP and FN analysis }

Feature Space Overlap (Class 3 $\rightarrow$ Class 2). Consistent
with our baseline observations, the models systematically misclassify
Class 3 as Class 2. To investigate the root cause, we isolated the
misclassified instances and compared their feature profiles against
correctly identified Class 2 samples (we used specifically Decision
Tree model predictions to choose the predicted labels).\\
we observed that for the followinng features in Class 3 were indistinguishable
from Class 2:\\
'inbound', 'dport\_443', 'dport\_80', 'dport\_53', 'dport\_22', 'dport\_0',
'protocol\_17', 'dport\_465', 'src\_port\_medium', 'protocol\_6','src\_port\_high'\\
and that the top 5 diffrences were (feature name and average distance
respectively ):\\
PC2 0.003494, PC1 0.003982, PC6 0.004872, PC5 0.005935, PC16 0.007088\\
\\
This error technically represents a False Negative for Class 3 (we
failed to identify the specific attack type) and a False Positive
for Class 2 (we raised a specific alarm for the wrong attack).\\
This error might not be critical due to the fact that we are confusing
between attacks and not attacks and benign flows. In case those attacks
require very different handling by the network security this error
will be more critical 

\subsubsection{Real-World Implications (FP, FN, Recall, Precision)}

assuming true is recognition of any ddos attack\\
\textbf{FP }-\textbf{ }The \textquotedbl False Alarm\textquotedbl{}
A False Positive occurs when benign traffic is incorrectly blocked
as malicious. This harms user experience and causes direct revenue
loss. We consider this the worst-case scenario, as it creates a \textquotedbl Self-Denial
of Service\textquotedbl{} where the defense system itself prevents
legitimate access.\\
\textbf{FN} - The \textquotedbl Missed Detection\textquotedbl{} A
False Negative occurs when attack traffic mistakenly enters the network.
While ideally avoided, the network can ``absorve'' those attacks
and keep working regularly. Therefore, FNs are more tolerable than
FPs in our specific context.\\
\textbf{Precision }- Optimizing for Precision ensures that when the
model raises an alarm, we can be confident it is a genuine threat,
minimizing the risk of blocking legitimate traffic.\\
\textbf{Recall -} This metric penalizes False Negatives. While important,
prioritizing Recall too heavily would likely lead to an over-aggressive
model that detects more attacks but also \textquotedbl throws away\textquotedbl{}
valid benign data. Given our priority to protect user experience,
we accept a lower Recall to guarantee higher Precision.

\subsubsection{split comparison}

as we saw throughout, this section our model didn't have significant
differences in the two splits (beside one regarding miss classification
one of the attacks only on GNB baseline model which was elaborated
earlier) 



\section{Unsupervised learning -- clustering}

\subsection{Experimental Setup}
We applied K-Means, Gaussian Mixture Models (GMM), and DBSCAN to the PCA-reduced features. Target labels were removed during training to ensure a strictly unsupervised process and used only for external validation (ARI/RI).

\subsection{Determining the number of clusters (k)}
We evaluated $k\in[2,16]$ using the \textbf{Elbow Method} (Inertia) to measure compactness and the \textbf{Silhouette Score} to measure separation. The optimal $k$ was selected where the Inertia score \textbf{flattened out} and the Silhouette Score \textbf{peaked}.
\paragraph{kmeans:}


\begin{figure}[H]
    \centering
    \begin{subfigure}[b]{0.48\linewidth}
        \includegraphics[width=\linewidth]{k_evaluation/kmeans_silhouette}
        \caption{Silhouette}
    \end{subfigure}
    \hfill
    \begin{subfigure}[b]{0.48\linewidth}
        \includegraphics[width=\linewidth]{k_evaluation/kmeans_inertia}
        \caption{Inertia}
    \end{subfigure}
    \caption{K-Means k Evaluation}
\end{figure}

based on the peak in the silhouette and steady declining in inertia
we chose k=8 for kmeans.

\paragraph{gmm:}

\begin{figure}[H]
    \centering
    \begin{subfigure}[b]{0.48\linewidth}
        \includegraphics[width=\linewidth]{k_evaluation/gmm_silhouette}
        \caption{Silhouette}
    \end{subfigure}
    \hfill
    \begin{subfigure}[b]{0.48\linewidth}
        \includegraphics[width=\linewidth]{k_evaluation/gmm_inertia}
        \caption{Inertia}
    \end{subfigure}
    \caption{GMM k Evaluation}
\end{figure}


based on the pick in the silhouette and steady declining in inertia
we chose k=13 for kmeans.

(dbscan does not require the number k in adcance)

\subsection{Evaluating ``Natural Grouping'' and Fine-Tuning}
To validate if the unsupervised clusters aligned with known attack families, we compared the clustering results against Ground Truth labels using \textbf{Rand Index (RI)} and \textbf{Adjusted Rand Index (ARI)}. We then performed a grid search to maximize the Silhouette Score and evaluated if better well-separated clusters corresponded to better attack detection (higher ARI).

\subsubsection{Hyperparameter Tuning + natural grouping}

To ensure robust results, we performed a grid search for each algorithm.
This involved testing combinations of parameters (e.g., distance metrics
like Euclidean vs. Cosine) to maximize the Silhouette Score. We then
evaluated the ARI/RI on these optimized models to see if \textquotedbl better\textquotedbl{}
clusters (geometrically) corresponded to \textquotedbl better\textquotedbl{}
attack detection.

\paragraph{kmeans}
\leavevmode

\begin{figure}[H]
    \centering
    \begin{subfigure}[b]{0.32\linewidth}
        \includegraphics[width=\linewidth]{k_evaluation/grid_kmeans_silhouette}
        \caption{Silhouette}
    \end{subfigure}
    \hfill
    \begin{subfigure}[b]{0.32\linewidth}
        \includegraphics[width=\linewidth]{k_evaluation/grid_kmeans_ri}
        \caption{RI}
    \end{subfigure}
    \hfill
    \begin{subfigure}[b]{0.32\linewidth}
        \includegraphics[width=\linewidth]{k_evaluation/grid_kmeans_ari}
        \caption{ARI}
    \end{subfigure}
    \caption{K-Means Grid Search Metrics}
\end{figure}

We selected the Elkan algorithm with k-means++ for its superior internal quality (Silhouette: 0.571), indicating distinct geometric structures. However, external validation revealed a significant divergence from ground truth (ARI: 0.144). Notably, random initialization yielded better label alignment (ARI $\approx$ 0.344). This inverse relationship suggests that while K-Means identifies strong structural groupings in the feature space, these geometric clusters do not correspond to the semantic attack categories defined in the labels.


\paragraph{gmm}
\leavevmode

\begin{figure}[H]
    \centering
    \begin{subfigure}[b]{0.32\linewidth}
        \includegraphics[width=\linewidth]{k_evaluation/grid_gmm_silhouette}
        \caption{Silhouette}
    \end{subfigure}
    \hfill
    \begin{subfigure}[b]{0.32\linewidth}
        \includegraphics[width=\linewidth]{k_evaluation/grid_gmm_ri}
        \caption{RI}
    \end{subfigure}
    \hfill
    \begin{subfigure}[b]{0.32\linewidth}
        \includegraphics[width=\linewidth]{k_evaluation/grid_gmm_ari}
        \caption{ARI}
    \end{subfigure}
    \caption{GMM Grid Search Metrics}
\end{figure}

We selected k-means initialization with diagonal covariance, which maximized internal quality (Silhouette: 0.545). Unlike K-Means, this model showed strong alignment with ground truth (ARI: 0.371, RI: 0.851), demonstrating a positive correlation between geometric fit and label accuracy. This significant improvement over K-Means (ARI: 0.144) suggests that the ``natural'' shape of attack traffic is better approximated by Gaussian distributions with independent variances than by spherical clusters



\paragraph{dbscan - euclidean}
\leavevmode

\begin{figure}[H]
    \centering
    \begin{subfigure}[b]{0.32\linewidth}
        \includegraphics[width=\linewidth]{k_evaluation/grid_euclidean_dbscan_silhouette}
        \caption{Silhouette}
    \end{subfigure}
    \hfill
    \begin{subfigure}[b]{0.32\linewidth}
        \includegraphics[width=\linewidth]{k_evaluation/grid_euclidean_dbscan_ri}
        \caption{RI}
    \end{subfigure}
    \hfill
    \begin{subfigure}[b]{0.32\linewidth}
        \includegraphics[width=\linewidth]{k_evaluation/grid_euclidean_dbscan_ari}
        \caption{ARI}
    \end{subfigure}
    \caption{DBSCAN (Euclidean) Grid Search Metrics}
\end{figure}

We selected eps=0.05 and min\_samples=20, achieving the highest Silhouette Score (0.888) and a strong Rand Index (0.902)1. This exceptionally high internal quality suggests the successful identification of dense ``core'' traffic patterns. While this configuration yielded an ARI of 0.332, relaxing the density constraint ($eps=0.5$) further improved alignment (ARI 0.383). These results confirm that attack families are better defined by density continuity than by distance from a centroid, making density-based clustering far more suitable than spherical partitioning for this dataset

\paragraph{dbscan - cosine}
\leavevmode

\begin{figure}[H]
    \centering
    \begin{subfigure}[b]{0.32\linewidth}
        \includegraphics[width=\linewidth]{k_evaluation/grid_cosine_dbscan_silhouette}
        \caption{Silhouette}
    \end{subfigure}
    \hfill
    \begin{subfigure}[b]{0.32\linewidth}
        \includegraphics[width=\linewidth]{k_evaluation/grid_cosine_dbscan_ri}
        \caption{RI}
    \end{subfigure}
    \hfill
    \begin{subfigure}[b]{0.32\linewidth}
        \includegraphics[width=\linewidth]{k_evaluation/grid_cosine_dbscan_ari}
        \caption{ARI}
    \end{subfigure}
    \caption{DBSCAN (Cosine) Grid Search Metrics}
\end{figure}


We observed a distinct performance split: using eps=0.3 achieved a remarkably high Silhouette Score (0.859) but suffered a critical failure in external validation (ARI: 0.000). This demonstrates a `collapse' phenomenon where, despite high geometric separation, the algorithm likely grouped the majority of flows into a single cluster or noise. While a lower threshold ($eps=0.05$) improved label alignment (ARI: 0.201), the resulting geometric structure was negligible (Silhouette $\sim$ 0.021), confirming that Cosine distance is ineffective for separating these specific attack types
\\
\textbf{Comparison} (Euclidean vs. Cosine): Euclidean distance significantly
outperformed Cosine distance for this specific dataset. it has better
Silhouette score going hand in hand with the good ARI indicating good
clustering in terms of good clustering separation and correlation
to our original labels

\subsubsection{visualization for best params}
\leavevmode

% Page 1: First two visualizations
\begin{figure}[htbp]
    \centering
    \begin{subfigure}[b]{0.45\linewidth}
        \includegraphics[width=\linewidth]{k_evaluation/kmeans_visualization}
        \caption{K-Means Visualization}
    \end{subfigure}
    \hfill
    \begin{subfigure}[b]{0.45\linewidth}
        \includegraphics[width=\linewidth]{k_evaluation/gmm_visualization}
        \caption{GMM Visualization}
    \end{subfigure}
    \caption{Clustering Visualizations: K-Means and GMM}
\end{figure}

% Page 2: Remaining two visualizations
\begin{figure}[htbp]
    \centering
    \begin{subfigure}[b]{0.45\linewidth}
        \includegraphics[width=\linewidth]{k_evaluation/dbscan_visualization}
        \caption{DBSCAN Visualization}
    \end{subfigure}
    \hfill
    \begin{subfigure}[b]{0.45\linewidth}
        \includegraphics[width=\linewidth]{k_evaluation/ground_truth_visualization}
        \caption{Ground Truth}
    \end{subfigure}
    \caption{Clustering Visualizations: DBSCAN and Ground Truth}
\end{figure}

\FloatBarrier
\subsubsection{coarse analysis of detected clusters}

Analysis ECDF of number of flows per cluster reveals distinct structural differences:

\begin{figure}[H]
    \centering
    \begin{subfigure}[b]{0.32\linewidth}
        \includegraphics[width=\linewidth]{k_evaluation/gmm_ecdf}
        \caption{GMM ECDF}
    \end{subfigure}
    \hfill
    \begin{subfigure}[b]{0.32\linewidth}
        \includegraphics[width=\linewidth]{k_evaluation/dbscan_ecdf}
        \caption{DBSCAN ECDF}
    \end{subfigure}
    \hfill
    \begin{subfigure}[b]{0.32\linewidth}
        \includegraphics[width=\linewidth]{k_evaluation/kmeans_ecdf}
        \caption{K-Means ECDF}
    \end{subfigure}
    \caption{ECDF of Flows per Cluster}
\end{figure}


\begin{itemize}
    \item \textbf{K-Means:} dominated by one massive cluster (>43k flows), 5 medium clusters (ranging between 2,300 and 8,700 flows) and 2 negligible clusters (under 100 flows) .
    \item \textbf{GMM:} produced a balanced distribution with 6 small(under 1000 flows), 5 medium (under 6000 above 1000), and 2 large clusters(between
17000 and 25000 flows), better reflecting the attack hierarchy.
    \item \textbf{DBSCAN:} fragmented the data into 110 clusters (mostly <1,000 flows), likely over-segmenting dominant traffic into micro-clusters and remove of noise.
\end{itemize}


\section{Clusters explainability and analysis}

This final section examines the internal structure of our generated clusters to identify behavioral patterns and evaluate their alignment with physical attack categories. 

\subsection{Model Selection and Experimental Setup}
We selected Gaussian Mixture Models (GMM) as it provides the best balance between statistical performance and structural alignment1111. Unlike DBSCAN, which over-segmented traffic into 110 micro-clusters, GMM's $k=13$ more closely matches the original attack classes, ensuring better interpretability2222. Additionally, GMM significantly outperformed K-Means in label alignment (ARI: 0.371 vs 0.144), proving that modeling complex Gaussian distributions is more effective for characterizing DDoS patterns than spherical models
\subsection{Cluster Distribution of Traffic Types}

To visualize the relationship between clusters and ground truth (GT) labels, we generated two normalized heatmaps and an Empirical Cumulative Distribution Function (ECDF).

\begin{figure}[H]
    \centering
    \begin{subfigure}{0.48\linewidth}
        \centering
        \includegraphics[width=\linewidth]{heatmap_row_normalized_pct.png}
        \caption{Row-normalized: Class concentration per cluster.}
    \end{subfigure}
    \hfill
    \begin{subfigure}{0.48\linewidth}
        \centering
        \includegraphics[width=\linewidth]{heatmap_col_normalized_pct.png}
        \caption{Column-normalized: Cluster purity by class.}
    \end{subfigure}
    \caption{GMM Cluster-Class Contingency Heatmaps}
\end{figure}
The row-normalized heatmap illustrates whether a specific class is concentrated or distributed across many clusters, while the column-normalized heatmap identifies which components (attacks) construct each cluster. 

\begin{figure}[H]
    \centering
    \includegraphics[width=0.6\linewidth]{ecdf_clusters_per_class.png}
    \caption{ECDF of number of clusters assigned to each class}
\end{figure}

The ECDF curve shows, for each k, the percentage of classes that are distributed in at most k clusters:
low values of k indicate well-concentrated classes (specific patterns),
high values indicate classes more fragmented across multiple clusters (greater heterogeneity or overlap with other classes)
Our ECDF highlights that many classes are relatively concentrated (most appear in fewer than 5–6 clusters),
but there are also very heterogeneous/dispersed classes that are present in a high number of clusters (up to 9).
This confirms what was observed in the heatmaps: some traffic types have a very specific signature and
are isolated into a few clusters, while others (typically benign traffic with more varied or overlapping patterns)
are distributed across more clusters\\

By examining the column-normalized heatmap, we identify several ``pure'' clusters where all elements belong to a single class, specifically clusters 3, 6, 9, and 12. We can deduce that these clusters capture highly specific traffic patterns. Furthermore, analysis of the absolute count contingency table reveals that cluster 9 contains only a single isolated element.\\

Based on these observations, it is evident that while the clusters reflect the GT labels, they do so only partially, as there is no global 1-to-1 correspondence. Although pure clusters exist, the classes corresponding to them are often distributed across multiple other clusters. Conversely, certain classes—such as \textit{ddos\_snmp}, \textit{ddos\_udp}, \textit{ddos\_udp\_lag}, and \textit{ddos\_netbios}—are localized primarily or entirely within a single cluster that is nonetheless mixed, containing multiple attack types. This confirms that some attacks possess traffic patterns so similar that they are grouped together in the same unsupervised cluster.\\

A part of benign traffic is clearly distinct from attacks and is well separable,
in fact clusters 6 and 12, which are pure, are 100% benign.
We note however that not all benign traffic falls into those clusters but is also distributed across others, which are heterogeneous
(8=16.6%, 4=16%, 6=13.4%, 1=12.4%, etc.)
This means that a part of benign traffic is grouped together with attacks and shares with them similar patterns.
This makes sense because benign traffic can be very heterogeneous
and many DDoS reflection/amplification attacks use protocols/ports that also exist in legitimate traffic
(e.g. DNS on port 53/UDP: legitimate DNS queries vs DNS amplification,
HTTP/HTTPS on ports 80/443: normal web traffic vs flood )

\subsection{Feature Importance Analysis}

We adopted a model-agnostic explainability approach using a \textit{RandomForestClassifier} to identify the primary drivers of cluster formation, achieving a reconstruction accuracy of 0.9986. Permutation Importance was computed on the test set using the \textit{f1\_macro} metric.

\begin{figure}[H]
    \centering
    \includegraphics[width=0.7\linewidth]{feature_importance_top15.png}
    \caption{Top 15 features identified via Permutation Importance}
\end{figure}

\begin{itemize}
    \item \textbf{Top Discriminators:} \texttt{dport\_80} is the most significant feature, indicating that traffic towards HTTP is a main separation driver. 
    \item \textbf{Directionality:} The \texttt{inbound} feature follows in importance, confirming that flow direction contributes markedly to distinguishing groups.
    \item \textbf{Protocol Influence:} \texttt{dport\_443} distinguishes HTTPS traffic, while \texttt{protocol\_17} (UDP) identifies specific patterns associated with reflection-based attacks.
\end{itemize}



\end{document}
